% Generated draft: requires mathematical and source review.
\documentclass[11pt]{article}
\usepackage[margin=25mm]{geometry}
\usepackage{fontspec}
\setmainfont{Latin Modern Roman}
\usepackage{amsmath,amssymb,mathtools,unicode-math}
\setmathfont{Latin Modern Math}
\usepackage{enumitem,needspace,xurl,hyperref}
\setlength{\parindent}{0pt}
\setlength{\parskip}{5pt}
\setlength{\emergencystretch}{3em}
\newcommand{\N}{\mathbb{N}}
\newcommand{\Z}{\mathbb{Z}}
\newcommand{\Q}{\mathbb{Q}}
\newcommand{\R}{\mathbb{R}}
\newcommand{\C}{\mathbb{C}}
\newcommand{\F}{\mathbb{F}}
\newcommand{\A}{\mathbb{A}}
\newcommand{\PP}{\mathbb{P}}
\newcommand{\E}{\mathbb{E}}
\newcommand{\eps}{\varepsilon}
\newcommand{\dd}{\,\mathrm{d}}
\begin{document}
\section*{UnsolvedMath ID: 30005849}
\textbf{Source problem code:} catalogue entry
\subsection{Definitions and mathematical statement}
The supplied record asks for a formulation of Geroch stability analogous to the good-set/bad-set positive-mass stability theorem attributed there to Dong and Song. It does not supply the hypotheses or conclusion of that theorem, or a complete definition of the intended Geroch stability statement. Accordingly, the following is a research target with explicitly unresolved parameters, not a fully specified conjecture.

For the Riemannian manifolds denoted in the excerpt by \((M_i,g_i)\), write \(Z_i\subseteq M_i\) for a proposed bad set and \(G_i=M_i\setminus Z_i\) for its complement. The restricted length distance mentioned in the source is
\[
\widehat d_{g_i}(x,y)=\inf\{L_{g_i}(\gamma):\gamma\text{ is a rectifiable curve contained in }G_i\text{ joining }x\text{ to }y\},\qquad x,y\in G_i.
\]
Here \(L_{g_i}\) is Riemannian curve length, and the infimum is \(+\infty\) if there is no such curve. Thus additional conditions are needed to ensure that this construction yields finite distances on the intended space. It need not equal the restriction of the ambient Riemannian distance.

\textbf{Requested research problem.} Formulate and prove a Geroch stability result with a good-set/bad-set decomposition analogous to the Dong--Song result referenced as ``Theorem 4'' in the supplied excerpt. Specify the admissible sequence, its scalar-curvature and other geometric hypotheses, the construction and quantitative smallness requirements for \(Z_i\), and the comparison or limit space. The excerpt mentions pointed measured Gromov--Hausdorff convergence for \(G_i\) equipped with \(\widehat d_{g_i}\), with \(p_i\in G_i\) \emph{any choice of base point}. If this is intended as the requested conclusion, its schematic form is
\[
(G_i,\widehat d_{g_i},\mu_i,p_i)\longrightarrow (X,d,\mu,p)
\quad\text{in the pointed measured Gromov--Hausdorff sense}.
\]
The measures \(\mu_i\), limit \((X,d,\mu,p)\), any completion or subsequence convention, and the precise scope of the base-point quantifier are not specified. The excerpt's disrupted syntax also leaves uncertain whether this convergence clause describes Theorem 4 or the requested Geroch analogue.

\textbf{Related questions retained from the excerpt.} The passage labelled Question 3 asks whether the bad sets \(Z_i\) exhibit only the phenomena called ``splines'', ``bubbles'', and ``drawstrings''. It proposes seeking one stability conjecture whose assumptions rule out splines and bubbles and whose conclusion is \(d_p\) stability, and another whose assumptions rule out bubbles and drawstrings and whose conclusion is volume-preserving Sormani--Wenger intrinsic flat convergence. These terms, the parameter \(p\), and the exclusion conditions are not defined in the record. The passage presents the three prospective stability results as evidence for this proposed description of bad sets, not as an established classification. It additionally qualifies the discussion by observing that ``sewing'' is possible when the topology of the sequence is allowed to vary.
\subsection{Short English statement}
The problem asks for a Geroch stability theorem that separates a sequence into good regions and controlled bad regions, following a result attributed to Dong and Song. The supplied excerpt mentions convergence using lengths of curves confined to the good regions and asks related questions about possible bad-region phenomena. Essential hypotheses, quantitative controls, convergence data, and several technical terms remain unspecified, so a complete conjecture requires mathematical review of the cited report.
\subsection{Sources}
\begin{itemize}
\item Local UnsolvedMath registry: supplied complete record for permanent ID \(30005849\), entitled \emph{Geroch Stability with a Good-Set Bad-Set Decomposition}; exact source problem code \texttt{catalogue entry}, proposed year 2024. No individual proposer is identified. The record attributes the motivating positive-mass stability theorem to Dong and Song but supplies no separate bibliographic citation for it.
\item \emph{Analysis, Geometry and Topology of Positive Scalar Curvature Metrics} (2024), Oberwolfach Reports, pp.~483--566, DOI \texttt{10.4171/owr/2024/9}: \url{https://doi.org/10.4171/owr/2024/9}. The supplied background describes Oberwolfach Reports as documenting mathematical workshops and their open-problem discussions. This is the supplied report-level citation; the relevant theorem text and a precise problem-page locator are absent.
\item Supplied literature source: \emph{Intrinsic-flat stability with boundary}, \url{https://eurekamag.com/research/096/975/096975145.php}.
\item Supplied literature source: \emph{Conjectures on convergence and scalar curvature}, \url{https://arxiv.org/abs/2103.10093}.
\item Literature assessment dated exactly \texttt{2026-08-22}: the supplied curation labels the problem partially solved and reports that versions of Geroch or VADB stability are proved for fixed manifolds and related weak-convergence hypotheses, while the general formulation requested in the report is not established. This is supplied dated curation derived from the OWR record, not an independently verified literature assessment. The supplied background expressly requests independent verification of time-sensitive status and cited claims.
\item The registry labels the statement \texttt{corrected\_verified} and says it was reconstructed and checked against the official Oberwolfach report. These are supplied provenance claims; no independent source visit or verification is asserted here.
\end{itemize}
\end{document}
